\begin{abstract}
    Recent work has explored the potential of graph neural networks (GNNs) for tackling combinatorial problems on graphs. This has motivated a systematic characterization of GNN expressivity, culminating (among other results) in a hierarchy of higher‑order architectures with provably increasing expressive power. In this work, we examine GNN expressivity via the \emph{clique--coloring problem} (CPP), a promise problem distinguishing graphs with a $k$-clique from those that are $(k\!-\!1)$-colorable. As a propositional theory, CPP is a canonical source of hard tautologies in proof complexity, while algorithmically CPP is famously polynomial‑time solvable via the computation of a graph's Lovász number.

    We prove that standard message-passing GNNs are fundamentally unable to distinguish these two classes of graphs. For every $k \ge 3$, we construct an infinite family of graph pairs, one containing a $k$-clique and the other $(k\!-\!1)$-colorable, that are indistinguishable by any such GNN. This result is unconditional and does not rely on complexity-theoretic assumptions. We show empirically that our hard graph family is near perfectly separated by higher order GNNs of order 3. Finally, we complement our theoretical findings with experiments demonstrating that while CPP is easy in general for GNNs on planted random graph distributions, performance deteriorates when the planted $k$-clique and $(k\!-\!1)$-colorable graph distributions are adjusted to have matching average degree.

\end{abstract}